\section{Introduction}
Large language models (LLMs) demonstrates a level of reasoning~\citep{qiao2023reasoning} but falter on multi-step reasoning unless given a clear plan or scaffold~\citep{wei2022chain}. They suffer from the hallucination problem, where they produce plausible-sounding yet incorrect output~\citep{huang2025survey}, and can be distracted by irrelevant parts of the input \citep{freda2023distracted}. To improve multi-step reasoning, prompt engineering (Definition~\ref{def:prompt_engineering_problem}) builds \emph{prompt schemes} that pair fixed instruction text with an explicit algorithmic procedure \citep{liu2023pre}. These schemes define a small, principled set of actions for the model to execute over several intermediate steps—often called \emph{thoughts}, i.e., discrete units in a multi-step solution \citep{wei2022chain, besta2024graph, besta2024demystifying}. Consequently, recent studies concentrate on mathematical and logical reasoning tasks \citep{zhou2023leasttomost, shin2025enhancing}. Success in these settings shows that, under well-designed prompts, LLMs can plan, execute, and check intermediate steps rather than jumping to a final guess, leading to more reliable solutions.

However, previous works often overlook the \textit{scalability issue} in prompt engineering. As explained in Definition~\ref{def:scalability_problem}, scalability concerns whether a prompt scheme remains effective as the input size increases~\citep{chen2025survey}.\footnote{While recent works aim to increase the context length~\citep{chen2023extending, chen2024clex, ding2024longrope} to improve LLM performances over long inputs~\citep{levy2024same}, a larger context length does not inherently guarantee effective multi-step reasoning~\citep{liu2024lost}.} Existing methods generally fall into two categories. Monolithic prompting methods---e.g., Chain of Thought (CoT)~\citep{wei2022chain} and Least to Most (L2M)~\citep{zhou2023leasttomost}---perform a single~\citep{wei2022chain, kojima2022large, wang2023plan} or fix number~\citep{wang2022self, zhou2024self, dua2022successive, zhou2023leasttomost} of inference passes. Such designs usually scale poorly, as they lack modularity, i.e., the separation of reasoning steps into distinct inference passes, and are limited to chain-like structures that force information to accumulate in later steps, increasing the risk of error propagation~\citep{besta2024demystifying}. Algorithmic prompting methods---e.g., Tree of Thoughts (ToT)~\citep{yao2024tree} and Graph of Thoughts (GoT)~\citep{besta2024graph}---in contrast, organize reasoning as trees
% ~\citep{yao2024tree, sel2024algorithm} 
or graphs,
% ~\citep{besta2024graph, zebaze2024tree}
coordinating multiple operations through external control scripts, i.e., Python modules. While these methods improve modularity by separating reasoning steps into separate LLM inferences, they also introduce and rely on rigid, hard-coded manual workflows that make it difficult to adapt to varying instance sizes. Consequently, both categories struggle to scale reliably as task size increase.

To address the scalability issue, it is important for a prompt scheme to demonstrate two key capabilities: 1) \textit{step-length generalization}---the ability to scale the number of reasoning steps while keeping each step \textit{elementary}, i.e., the minimal reasoning units into which the problem can be decomposed, to effectively reduce error~\citep{zhou2023leasttomost, xue2024decompose} (see the theoretical analysis in Section~\ref{sec:theory}); and 2) \textit{step-dependency generalization}---the ability to maintain correct dependencies between steps and coordinate them through network-structured reasoning, correctly recombining the decomposed steps into a final answer as the instance size increases~\citep{zhou2024self, yao2024tree, besta2024graph}. Real-world tasks vary in size and complexity, requiring scalable prompt schemes~\citep{ji2023exploring}. For example, in code refactoring, a programmer may efficiently update a small script by modifying multiple components at once. However, for a larger codebase, breaking it into steps that process one component at a time reduces errors, but the coordination complexity between components usually grows beyond linear chains or trees as the system scales.

\begin{figure*}[t]
    \centering
    \includegraphics[width=\textwidth]{figs/visual-compare_xnot.jpg}
    % \vspace{-5mm}
    \caption{Prompt scheme comparison. While some monolithic prompting methods introduce modular execution (e.g., L2M), they remain limited to linear, chain-like structures. Algorithmic prompting supports networked reasoning but relies on rigid, hard-coded scripts. In contrast, \method dynamically generates a \textit{network of thought} without requiring predefined workflows or complex control scripts.}
    % \vspace{-7mm}
    \label{fig:vis-compare}
\end{figure*}
 % 

 % generate a tailored solution structure for each task instance.

In this paper, we introduce \textit{\methodfull (\method)}, a prompt scheme that addresses the scalability issue in LLM reasoning through \textit{intelligence amplification}. Our idea is to transform prompt engineering from manually hard-coding complete workflows to providing the minimal scaffolding that LLMs can autonomously expand. While approaches like ToT~\citep{yao2024tree} and GoT~\citep{besta2024graph} require humans to predefine rigid structures for each task size, \method shifts this burden to LLMs by enabling it to infer both task decomposition and step coordination from simple examples. This allows step-length generalization, as LLM dynamically determines an appropriate number of elementary steps without requiring humans to redesign scripts for different problem sizes. As input complexity increases, \method enables LLMs to naturally expand its reasoning network with additional elementary steps to match the growing demands of the task.,

At the core of \method is the \textit{\scriptfull (\script)}, a structured text format that enables step-dependency generalization through a simple notation system. \script allows each reasoning step to reference outputs from any previous step, forming flexible networks of thought with precisely defined dependencies. Unlike the rigid chains~\citep{wei2022chain}, trees~\citep{yao2024tree}, or split-then-merge graphs~\citep{besta2024graph} used in prior work, \script empowers \method to construct instance-specific reasoning networks that capture the exact dependencies required for each query. By combining elementary step decomposition with dynamic dependency management, \method achieves \textit{intelligence amplification}: LLM generalizes the user-specified minimal examples to autonomously build comprehensive solution structures for previously unseen problems, eliminating the need for hard-coded workflows of each task variation.

Experiments demonstrate that \method scales much more reliably to larger input sizes. For example, when the sorting task scales from $16$ to $64$ numbers, most baselines drop to $0\%$ accuracy, whereas \method maintains $27\%$. \method also significantly reduces manual prompt engineering effort (up to $84.4\%$ less than ToT and $87.3\%$ less than GoT) and lowers the API costs compared to sophisticated algorithmic prompting methods. In addition, we provide theoretical analyses showing how decomposition into elementary steps improves performance under super-linear accuracy decay, while dynamic merge strategies reduce computational overhead compared to fixed structures. 
% Further ablation studies and an investigation into scaling from the 4-number to the 5-number variant of the Game of 24~\citep{yao2024tree} are presented in Appendix~\ref{app:experiment}.



% We validate \method across five scalable benchmarks that allows clear divisions by instance size over three domains: natural language (keyword counting), symbolic operations (sorting, set intersection), and numeric reasoning (arithmetic calculation, large-digit addition). We select these benchmarks to measure performance degradation as length increases. 

% It is worth noting that LLMs are inherently limited by a fixed context length. As input grows, performance degrades~\citep{liu2024lost, levy2024same}, since positional embeddings support only a bounded number of tokens~\citep{chen2023extending}. Although recent advances have extended context length~\citep{chen2023extending, chen2024clex, ding2024longrope}, they are unlikely to keep pace with the increasing complexity of real-world tasks. This highlights the need for prompt schemes that scale effectively with both input size and structure.
